% !PRINTED 291
\weberpage{291}

Zu diesen Invarianten geh\"ort auch der Grad $n$ der Gruppe $P$ selbst, der durch
die $j$ so bestimmt ist:
\weq{2}{n = j_1 j_2 \cdots j_{\mu-1} j_\mu.}
Denn nach der Bedeutung der Indices ist $n : j_1$ der Grad von $P_1$, $n : j_1 j_2$
der Grad von $P_2$ u.\,s.\,f., und da der letzte der Grade gleich $1$ ist, so ergiebt
sich die Formel (2).

Im Allgemeinen ist in einer Compositionsreihe von $P$ jedes Glied Normaltheiler nur
des n\"achst vorangehenden. Es k\"onnen aber auch einzelne Glieder vorkommen, die
auch noch von weiter vorangehenden Gliedern Normaltheiler sind. Besonders wichtig
sind solche Glieder, die Normaltheiler von $P$ selbst und also auch von allen ihnen
vorangehenden Gliedern der Compositionsreihe sind. Wir wollen sehen, welche
algebraische Consequenzen aus diesem Umstande zu ziehen sind.

Es sei $P_\nu$ ein Glied der Reihe (1), welches zugleich Normaltheiler von $P$ ist.
Der Index von $P_\nu$ in Bezug auf $P$ ist $j_1 j_2 \ldots j_\nu$, und dies Product
ist der Grad der Partialresolvente, durch die die Gruppe $P$ auf $P_\nu$ reducirt
wird, die wir mit $\chi(y) = 0$ bezeichnen wollen. Die Gruppe dieser Resolvente, die
eine Normalgleichung ist, erhalten wir nach dem Satze 1. in der Form $P/P_\nu$.

Wir k\"onnen nun leicht eine Compositionsreihe f\"ur diese Gruppe nebst der
zugeh\"origen Indexreihe finden, n\"amlich:
\weq{3}{
\begin{gathered}
P/P_\nu,\ P_1/P_\nu,\ P_2/P_\nu,\ \ldots,\ P_{\nu-1}/P_\nu,\ 1 \\
j_1,\quad j_2,\quad \ldots,\quad j_{\nu-1},\quad j_\nu.
\end{gathered}
}
Denn nach dem Satze 1., \S.\,6 ist $P_1/P_\nu$ ein Normaltheiler von $P/P_\nu$ vom
Index $j_1$, und es ist ein gr\"osster Normaltheiler,
weil nach 2., \S.\,6 \"uber $P_1/P_\nu$ kein Normaltheiler von $P/P_\nu$ stehen
kann, wenn \"uber $P_1$ kein Normaltheiler von $P$ steht. Und ebenso kann man in
Bezug auf die folgenden Glieder von (3) schliessen.

Daraus ziehen wir noch eine wichtige Folgerung. Wir haben schon im \S.\,7 gezeigt,
dass sich, wenn $Q$ irgend ein Normaltheiler von $P$ ist, eine Compositionsreihe von
$P$ finden l\"asst, in der $Q$ vorkommt. Ist nun $R$ ein anderer Normaltheiler von
$P$ und zugleich ein Theiler von $Q$, so kann man die Compositionsreihe von $P$ so
einrichten, dass $Q$ und $R$ darin vorkommen.

Wir denken uns eine solche Compositionsreihe bestimmt:
\weq{4}{P,\ Q',\ Q'',\ \ldots,\ Q,\ \ldots,\ R.}
